\chapter{Metriche di valutazione}
\label{cap:metriche}

Le metriche non sono uno strumento neutrale di misura: sono parte del disegno
sperimentale, e su un corpus saturo come quello descritto nel
Capitolo~\ref{cap:dataset} la scelta della metrica determina la conclusione. Questo
capitolo documenta la loro implementazione, i loro difetti misurati, e le
metriche introdotte per rimediarvi.

L'implementazione è in \texttt{src/utils/metrics.py}. Il modulo espone una
costante \texttt{METRICS\_VERSION}, il cui valore corrente è $2$: tutti i numeri
riportati in questa relazione sono calcolati con la versione $2$, e la costante
esiste perché confrontare numeri prodotti da versioni diverse del codice di
valutazione è un errore facile da commettere e difficile da individuare.

\section{Il difetto di ROUGE-L, misurato}
\label{sec:difetto-rouge}

L'implementazione di ROUGE-L impiegata (che segue la convenzione più diffusa
nella letteratura di riferimento, così da restare confrontabile con i numeri
pubblicati) esegue due normalizzazioni prima del confronto:

\begin{enumerate}
  \item porta tutto in minuscolo (\emph{case-insensitive});
  \item spezza i token sui caratteri non alfanumerici.
\end{enumerate}

Su un compito generico sono scelte innocue. Su un compito in cui il gloss è
distinto dal testo \emph{proprio} dal maiuscolo e dai prefissi con trattino, sono
devastanti. Ecco le misure dirette.

\subsection{Il maiuscolo non conta}

\begin{center}
\begin{tabular}{llr}
\toprule
Generato & Riferimento & ROUGE-L \\
\midrule
\texttt{cat sat} & \texttt{CAT SAT} & $1{,}00$ \\
\texttt{the cat sat} & \texttt{CAT SAT} & $0{,}80$ \\
\bottomrule
\end{tabular}
\end{center}

La prima riga è il problema in forma pura: un'uscita in minuscolo, cioè inglese
non convertito in gloss, ottiene il punteggio massimo. La seconda mostra che
anche mantenendo un articolo che il gloss elimina si resta a $0{,}80$.

Poiché il modello base, non addestrato, produce esattamente inglese in
minuscolo, ne segue che ROUGE-L assegna un punteggio alto a un sistema che non
ha eseguito \emph{nessuna} conversione. È la spiegazione del risultato
apparentemente paradossale della sezione~\ref{sec:trie-risultato}: il modello
senza vincolo ottiene $0{,}3648$ facendo la cosa sbagliata.

\subsection{I prefissi grammaticali vengono distrutti}

La tokenizzazione sui caratteri non alfanumerici spezza i marcatori del gloss:

\begin{center}
\begin{tabular}{ll}
\toprule
Stringa & Token prodotti \\
\midrule
\texttt{DESC-GOOD X-WE} & \texttt{['desc', 'good', 'x', 'we']} \\
\bottomrule
\end{tabular}
\end{center}

Il prefisso \texttt{DESC-}, che porta informazione grammaticale, diventa un
token indipendente. La conseguenza è che due gloss diversi condividono metà dei
loro token:

\begin{center}
\begin{tabular}{llr}
\toprule
Generato & Riferimento & ROUGE-L \\
\midrule
\texttt{DESC-GOOD} & \texttt{DESC-BAD} & $0{,}50$ \\
\texttt{X-WE} & \texttt{X-I} & $0{,}50$ \\
\bottomrule
\end{tabular}
\end{center}

Confondere \emph{buono} con \emph{cattivo} costa metà del punteggio. Confondere
\emph{noi} con \emph{io} costa metà del punteggio. Un sistema che indovina il
prefisso e sbaglia sistematicamente la parola ottiene $0{,}50$ su ogni gloss
marcato.

\subsection{La conclusione operativa}

ROUGE-L, su questo corpus e con questa implementazione, è una metrica di
\emph{copia}: premia la sovrapposizione superficiale di caratteri alfanumerici
in minuscolo, che è precisamente la relazione che il corpus sintetico garantisce
per costruzione (sezione~\ref{sec:sintetico}).

Non è stata rimossa, per due ragioni: è l'unica metrica su cui esistano numeri
pubblicati confrontabili, e la sua degenerazione è essa stessa un risultato da
documentare. È stata invece \textbf{degradata} da metrica primaria a metrica di
confronto storico, e nel codice le metriche di questa famiglia sono raccolte
sotto la costante \texttt{SATURATED\_OVERLAP\_METRICS} così che ogni consumatore
sappia come trattarle.

\section{Le metriche primarie}

Al posto di ROUGE-L, le metriche promosse a primarie (costante
\texttt{PRIMARY\_METRICS}) sono due.

\subsection{Exact match normalizzato}

Confronto di uguaglianza dopo normalizzazione degli spazi. È binario, severo e
non manipolabile: non esiste una strategia che aumenti l'exact match senza
produrre l'uscita corretta.

Il suo pregio principale su questo corpus è che \emph{ha ancora risoluzione}.
Come mostrato nel Capitolo~\ref{cap:sft}, fra baseline a regole e SFT ROUGE-L
distingue $0{,}0067$ mentre l'exact match distingue $0{,}2265$: due ordini di
grandezza di differenza nella capacità discriminante.

\subsection{Non-copy-token accuracy}

È la metrica progettata specificamente per questo corpus, e il contributo
metodologico principale del capitolo.

L'idea: si considerano soltanto le posizioni in cui il gloss di riferimento
\emph{differisce} dal token sorgente maiuscolizzato. Su quelle posizioni, e solo
su quelle, si misura l'accuratezza. Formalmente, detto
$\mathcal{N} = \{ t : r_t \neq \mathrm{upper}(x_t) \}$ l'insieme delle posizioni
non banali,
\begin{equation}
  \label{eq:noncopy}
  \texttt{non\_copy\_token\_accuracy}
  \;=\;
  \frac{|\{ t \in \mathcal{N} : g_t = r_t \}|}{|\mathcal{N}|}.
\end{equation}

Sui $2\,000$ prompt di valutazione, $|\mathcal{N}| = 9\,350$ token: un campione
ampio, non un residuo marginale.

Il potere discriminante è evidente:

\begin{center}
\begin{tabular}{lrr}
\toprule
Sistema & ROUGE-L & Non-copy accuracy \\
\midrule
Base zero-shot senza Trie & $0{,}3648$ & $0{,}0006$ \\
Base zero-shot con Trie & $0{,}1373$ & $0{,}0432$ \\
GRPO da modello base & $0{,}6076$ & $0{,}4641$ \\
Baseline a regole & $0{,}9685$ & $0{,}9424$ \\
SFT & $0{,}9752$ & $0{,}9660$ \\
\bottomrule
\end{tabular}
\end{center}

La colonna di destra ordina i sistemi in modo interpretabile e, soprattutto,
assegna al modello base senza vincolo il valore che merita: $0{,}0006$, cioè
niente. La colonna di sinistra gli assegna $0{,}3648$, che è più del doppio di
quanto ottenga lo stesso modello con il vincolo di decodifica attivo.

\section{Validity: la definizione, versione 2}

La \emph{validity} misura la frazione di uscite ben formate. La definizione è
stata rivista dopo aver osservato che la versione iniziale accettava uscite
degeneri.

Nella versione $2$ un'uscita è dichiarata \textbf{non valida} se si verifica
almeno una delle condizioni:

\begin{enumerate}
  \item meno del $50\%$ dei suoi token appartiene al vocabolario gloss;
  \item ha più di $4$ token e il rapporto fra token distinti e token totali
        (\emph{unicità}) è inferiore a $0{,}3$.
\end{enumerate}

La prima condizione cattura le uscite che non sono gloss. La seconda cattura le
ripetizioni degeneri del tipo \texttt{HOUSE HOUSE HOUSE HOUSE HOUSE}, che
soddisfano la prima condizione perfettamente (tutti i token sono in
vocabolario) e sono comunque prive di valore. La soglia di $4$ token evita di
penalizzare uscite legittimamente brevi in cui una ripetizione può essere reale.

\section{Pass@$k$}

L'implementazione segue lo stimatore combinatorio non
distorto~\eqref{eq:passk}, con soglia di correttezza posta a ROUGE-L $\ge 0{,}3$.

La soglia va dichiarata perché il valore di pass@$k$ dipende interamente da essa
e $0{,}3$ è, su questo corpus, una soglia \emph{permissiva}: dato il difetto
documentato nella sezione~\ref{sec:difetto-rouge}, un'uscita che supera $0{,}3$
non è necessariamente una traduzione utile. Pass@$k$ va quindi letta come misura
di \emph{copertura} del campionamento, non di qualità.

\section{La metrica composta e la sua correzione}

Una metrica composta comoda è il prodotto fra qualità media e frazione di uscite
valide, che penalizza i sistemi che ottengono punteggi alti su poche uscite ben
formate:
\begin{equation}
  \label{eq:valid-rouge-v2}
  \texttt{valid\_rouge\_l\_mean}
  \;=\;
  \overline{\texttt{rouge}} \times \texttt{validity\_rate}.
\end{equation}

Questa forma ha però un difetto: il prodotto di due medie \emph{non} è la media
del prodotto. Se le uscite non valide fossero sistematicamente quelle con ROUGE-L
più alto (caso possibile: un'uscita che copia l'inglese può avere ROUGE-L alto e
validity nulla), la~\eqref{eq:valid-rouge-v2} sovrastimerebbe.

È stata quindi affiancata una versione per riga:
\begin{equation}
  \label{eq:valid-rouge-v3}
  \texttt{valid\_rouge\_l\_mean\_v3}
  \;=\;
  \frac{1}{N}\sum_{i=1}^{N}
  \begin{cases}
    \texttt{rouge}_i & \text{se l'uscita } i \text{ è valida},\\
    0 & \text{altrimenti.}
  \end{cases}
\end{equation}

Il confronto fra le due:

\begin{center}
\begin{tabular}{lrr}
\toprule
Cella & Versione 2 & Versione 3 \\
\midrule
Base few-shot & $0{,}4581$ & $0{,}4645$ \\
Base zero-shot con Trie & $0{,}1243$ & $0{,}1335$ \\
\bottomrule
\end{tabular}
\end{center}

Le differenze sono piccole ($+0{,}006$ e $+0{,}009$) e nella direzione opposta a
quella temuta: la versione per riga dà valori leggermente \emph{più alti}, il che
significa che le uscite non valide tendevano ad avere ROUGE-L sotto la media.
Entrambe le versioni sono riportate, così che il lettore possa verificare che la
scelta fra le due non cambia le conclusioni.

\section{Incertezza: bootstrap clusterizzato}

Tutti gli intervalli di confidenza riportati in questa relazione sono ottenuti
per bootstrap, con una precisazione necessaria: il ricampionamento è
\textbf{clusterizzato per prompt}.

La ragione è che la valutazione genera $N=5$ uscite per ogni prompt. Le cinque
uscite dello stesso prompt non sono osservazioni indipendenti: condividono la
frase di ingresso, la sua difficoltà e il suo riferimento. Ricampionare le
singole uscite come se fossero indipendenti gonfierebbe la dimensione effettiva
del campione di un fattore fino a cinque, restituendo intervalli di confidenza
artificialmente stretti.

Il bootstrap clusterizzato ricampiona i \emph{prompt} con reimmissione,
portandosi dietro tutte le uscite di ciascun prompt selezionato. È la procedura
corretta in presenza di questa struttura a gruppi, e produce intervalli più
larghi e onesti. Gli intervalli riportati nella sezione sul confronto fra SFT e
baseline a regole (Capitolo~\ref{cap:sft}) sono calcolati in questo modo.

\section{Il protocollo di valutazione, riassunto}

Perché i numeri del capitolo successivo siano interpretabili, il protocollo è
fissato e identico per tutte le celle:

\begin{center}
\begin{tabular}{ll}
\toprule
Parametro & Valore \\
\midrule
Prompt valutati & $2\,000$ (dal test set) \\
Generazioni per prompt & $N = 5$ \\
Strategia di decodifica & campionamento, $\tau = 0{,}7$ \\
Versione delle metriche & $2$ \\
Intervalli di confidenza & bootstrap clusterizzato per prompt \\
Riferimento obbligatorio & baseline a regole \\
\bottomrule
\end{tabular}
\end{center}

Il campionamento a temperatura $0{,}7$ invece del greedy è una scelta che va
motivata: consente di stimare pass@$k$ e di misurare la varianza delle uscite,
che è l'informazione necessaria per interpretare il supporto dei rollout
(sezione~\ref{sec:supporto}). Il costo è che i numeri di exact match sono
leggermente inferiori a quelli che si otterrebbero in greedy; poiché il confronto
fra celle usa lo stesso protocollo, i confronti restano validi.
